% ============================================================
%  PHU LUC
% ============================================================
\unnumberedchapter{PHỤ LỤC}

\unnumberedsection{Phụ lục A. Hướng dẫn cài đặt và chạy hệ thống}

Hệ thống dùng EF Core theo hướng Code-First nên cơ sở dữ liệu được dựng tự động
từ migrations, \textbf{không cần nhập tệp sao lưu} \code{.bak}. Cùng một mã nguồn
chạy giống hệt nhau trên macOS và Windows; điểm khác biệt duy nhất là \textbf{chuỗi
kết nối}, vốn nằm trong cấu hình chứ không nằm trong mã nguồn.

\unnumberedsubsection{A.1. Chuẩn bị công cụ}

\begin{codeblock}
dotnet --version                     # can tu 10.0 tro len
dotnet tool install --global dotnet-ef    # neu chua co
\end{codeblock}

\unnumberedsubsection{A.2. Dựng SQL Server}

Chọn một trong ba cách, kết quả như nhau.

\textbf{Cách A, Docker (khuyến nghị, giống nhau trên macOS và Windows).} macOS
dùng OrbStack hoặc Docker Desktop; Windows dùng Docker Desktop. Cùng một tệp
\code{docker-compose.yml}:

\begin{codeblock}
cd setup
docker compose up -d
# SQL Server lang nghe o localhost:1433
# tai khoan sa / CookAdvisor@2026
\end{codeblock}

Mật khẩu SA mặc định chỉ dùng cho môi trường phát triển cục bộ, có thể đổi bằng
biến môi trường \code{MSSQL\_\allowbreak{}SA\_\allowbreak{}PASSWORD} trước khi chạy lệnh trên. Trên máy
Apple Silicon, ảnh container chạy dưới giả lập amd64, cấu hình đã có sẵn trong
tệp compose.

\textbf{Cách B, SQL Server cài trực tiếp trên Windows:} cài bản Express hoặc
Developer, bật chế độ xác thực hỗn hợp. \textbf{Cách C, LocalDB (Windows, nhẹ
nhất cho demo):} có sẵn khi cài Visual Studio.

\unnumberedsubsection{A.3. Cấu hình chuỗi kết nối}

Chuỗi kết nối chứa mật khẩu nên \textbf{không được đưa vào mã nguồn}. Đặt qua
User Secrets:

\begin{codeblock}
cd src/CookingAdvisor
dotnet user-secrets set \
  "ConnectionStrings:DefaultConnection" "<chuoi phu hop>"
\end{codeblock}

\begin{table}[H]
\centering
\caption{Chuỗi kết nối theo từng môi trường}
\label{tab:chuoi-ket-noi}
\begin{tabularx}{\textwidth}{|P{4.0cm}|Y|}
\hline
\bfseries Môi trường & \bfseries Chuỗi kết nối \\ \hline
Docker (macOS hoặc Windows) &
  \codelong{Server=localhost,1433;Database=CookingAdvisor;User Id=sa;
  Password=CookAdvisor@2026;TrustServerCertificate=True} \\ \hline
SQL Server cài trực tiếp (Windows) &
  \codelong{Server=localhost;Database=CookingAdvisor;User Id=sa;
  Password=<mật khẩu của bạn>;TrustServerCertificate=True} \\ \hline
LocalDB (Windows) &
  \codelong{Server=(localdb)\textbackslash MSSQLLocalDB;Database=CookingAdvisor;
  Trusted\_\allowbreak{}Connection=True;TrustServerCertificate=True} \\ \hline
\end{tabularx}
\end{table}

Tham số \code{TrustServerCertificate=True} cần thiết vì chứng chỉ của môi trường
phát triển là chứng chỉ tự ký. Tiếng Việt được lưu bằng kiểu \code{nvarchar}
(Unicode) nên hiển thị đúng trên mọi môi trường.

\unnumberedsubsection{A.4. Tạo cơ sở dữ liệu và chạy ứng dụng}

\begin{codeblock}
# tu thu muc goc cua repository
# dung luoc do co so du lieu
dotnet ef database update --project src/CookingAdvisor
# chay ung dung
dotnet run --project src/CookingAdvisor
\end{codeblock}

Ở lần chạy đầu tiên, lớp \code{DbInitializer} nạp dữ liệu mẫu gồm danh mục,
nguyên liệu, 26 món ăn và tài khoản quản trị.

\unnumberedsubsection{A.5. Cách thay thế cho Windows: dùng script T-SQL}

Nếu máy Windows chưa cài công cụ dòng lệnh \code{dotnet-ef}, có thể dựng lược đồ
bằng script T-SQL đã xuất sẵn tại \codelong{setup/sql/InitialCreate.sql}. Mở tệp bằng
SQL Server Management Studio hoặc \code{sqlcmd}, kết nối tới máy chủ, chọn hoặc
tạo cơ sở dữ liệu \code{CookingAdvisor} rồi thực thi.

Script này có tính lũy đẳng: nó tự kiểm tra bảng \code{\_\allowbreak{}\_\allowbreak{}EFMigrationsHistory}
nên chạy lại nhiều lần không gây lỗi. Script được sinh ra từ toàn bộ chuỗi
migrations hiện có nên lược đồ giống hệt bản dựng qua EF Core; khi có migration
mới cần sinh lại script theo lệnh trong \code{setup/README.md}.

Lưu ý: script chỉ tạo \textbf{lược đồ} (bảng, cột, khóa ngoại), không chứa dữ liệu
mẫu. Sau khi thực thi xong vẫn cần chạy ứng dụng \textbf{một lần} để
\code{DbInitializer} nạp dữ liệu, vì phần nạp dữ liệu nằm trong mã C\# chứ không
nằm trong tệp SQL.

\unnumberedsubsection{A.6. Tài khoản mặc định}

\begin{table}[H]
\centering
\caption{Tài khoản quản trị mặc định của môi trường demo}
\label{tab:tai-khoan}
\begin{tabularx}{\textwidth}{|P{3.4cm}|Y|Y|}
\hline
\bfseries Vai trò & \bfseries Email & \bfseries Mật khẩu \\ \hline
Quản trị viên & \codelong{admin@cookingadvisor.local} & \code{Admin@2026!Cook} \\ \hline
\end{tabularx}
\end{table}

Tài khoản này được tạo tự động khi chạy ứng dụng lần đầu trên cơ sở dữ liệu
trống, chỉ dùng cho môi trường demo của đồ án.

\unnumberedsubsection{A.7. Chạy kiểm thử}

\begin{codeblock}
dotnet test src/CookingAdvisor.sln
\end{codeblock}

Bộ kiểm thử dùng nhà cung cấp EF Core In-Memory nên chạy được mà không cần khởi
động SQL Server.

\unnumberedsection{Phụ lục B. Cấu trúc mã nguồn}

\begin{codeblock}
CookingAdvisor/
  setup/                       Tai nguyen cai dat
    docker-compose.yml         Cau hinh container SQL Server
    sql/InitialCreate.sql      Script T-SQL dung luoc do
    README.md                  Huong dan cai dat
  src/
    CookingAdvisor.sln
    CookingAdvisor/            Project web
      Controllers/             Tang dieu khien
      Areas/Admin/             Khu vuc quan tri
      Services/                Tang nghiep vu, 3 thuat toan
      Models/                  Thuc the EF Core, kieu liet ke
      ViewModels/              Du lieu truyen cho View
      Data/                    AppDbContext va DbInitializer
      Views/                   Giao dien Razor
      TagHelpers/              Tag helper he thong bieu tuong
      Migrations/              Lich su thay doi luoc do
      wwwroot/                 CSS, JavaScript, phong chu, anh
    CookingAdvisor.Tests/      Du an kiem thu xUnit
  thesis/                      Tai lieu do an
    doc/                       Ban .docx cua bao cao
    pdf/                       Ban .pdf de nop
    latex/                     Nguon LaTeX ban bao cao chinh
    images/                    Anh chup giao dien
    diagrams/                  So do (ma nguon Mermaid va anh)
\end{codeblock}

\unnumberedsection{Phụ lục C. Nguồn ảnh minh họa món ăn}

Toàn bộ 26 ảnh món ăn sử dụng trong hệ thống được tải từ \textbf{Wikimedia
Commons}, đều thuộc các giấy phép cho phép sử dụng tự do (CC0, Public Domain, CC
BY, CC BY-SA). Ảnh đã được giảm kích thước về 900 px chiều ngang.

Danh sách đầy đủ gồm tên món, tên tệp, tác giả, giấy phép và đường dẫn tới trang
nguồn được lưu tại \codelong{src/CookingAdvisor/wwwroot/images/recipes/CREDITS.md}.
Bảng~\ref{tab:nguon-anh} trích một số mục tiêu biểu.

\begin{table}[H]
\centering
\caption{Một số nguồn ảnh minh họa món ăn tiêu biểu}
\label{tab:nguon-anh}
\begin{tabularx}{\textwidth}{|Y|Y|P{3.4cm}|}
\hline
\bfseries Món ăn & \bfseries Tác giả & \bfseries Giấy phép \\ \hline
Phở bò & Andy Li & CC0 \\ \hline
Bún chả & Phương Huy & Public Domain \\ \hline
Bún bò Huế & Baoothersks & CC BY-SA 4.0 \\ \hline
Cơm tấm thịt nướng & Tokeisan & Public Domain \\ \hline
Gỏi cuốn & Tran Hai Duong & CC0 \\ \hline
Chả giò & Satdeep Gill & CC BY-SA 4.0 \\ \hline
Bánh xèo & Orderinchaos & CC BY-SA 4.0 \\ \hline
Rau muống xào tỏi & Andy Wright & CC BY 2.0 \\ \hline
\end{tabularx}
\end{table}

\unnumberedsection{Phụ lục D. Mã nguồn tiêu biểu}

\unnumberedsubsection{D.1. Thuật toán gợi ý theo nguyên liệu}

{\raggedright Trích từ \codelong{src/CookingAdvisor/Services/SuggestionService.cs}:\par}

\begin{codeblock}
public async Task<List<RecipeSuggestionViewModel>>
    SuggestAsync(IReadOnlyCollection<int> ownedIngredientIds)
{
    var owned = ownedIngredientIds.ToHashSet();
    if (owned.Count == 0)
        return [];

    // Recipes sharing no ingredient with the owned set are
    // not worth suggesting.
    var recipes = await db.Recipes
        .Where(r => r.RecipeIngredients
            .Any(ri => owned.Contains(ri.IngredientId)))
        .Select(r => new { /* ... */ })
        .ToListAsync();

    return recipes
        .Select(r =>
        {
            var missing = r.Ingredients
                .Where(i => !owned.Contains(i.IngredientId))
                .Select(i => i.Name)
                .OrderBy(name => name)
                .ToList();

            return new RecipeSuggestionViewModel
            {
                TotalIngredientCount = r.Ingredients.Count,
                MatchedCount =
                    r.Ingredients.Count - missing.Count,
                MissingIngredients = missing
                /* ... */
            };
        })
        .OrderByDescending(s => s.CanCookNow)
        .ThenByDescending(s => s.Coverage)
        .ThenBy(s => s.MissingIngredients.Count)
        .ThenBy(s => s.Name)
        .ToList();
}
\end{codeblock}

\unnumberedsubsection{D.2. Xử lý trường hợp phải lặp món trong thuật toán lập
thực đơn}

{\raggedright Trích từ \codelong{src/CookingAdvisor/Services/MenuPlannerService.cs}. Đây là đoạn
khắc phục khiếm khuyết suy thoái tính đa dạng đã trình bày ở
mục~\ref{sec:khiem-khuyet}:\par}

\begin{codeblock}
var suitable = candidates
    .Where(c => c.SuitableMealTypes.HasFlag(flag)).ToList();
var pool = suitable.Where(c => !used.Contains(c.Id)).ToList();

// Repeats only become necessary once every suitable dish
// has been used; meal-type suitability is relaxed after that.
var repeating = pool.Count == 0;
if (repeating) pool = suitable;
if (pool.Count == 0) pool = candidates;

var ordered = repeating
    // Variety beats preference once repeats are unavoidable:
    // without this key the calorie tie-break returns the same
    // dish for every remaining slot, so one dish fills the
    // week and the rest go unused.
    ? pool.OrderBy(c => useCounts.GetValueOrDefault(c.Id))
        .ThenByDescending(c => favoriteIds.Contains(c.Id))
    : pool.OrderByDescending(c => favoriteIds.Contains(c.Id));

var chosen = ordered
    .ThenBy(c => Math.Abs(
        dayCalories + c.CaloriesPerServing - mealTarget))
    .ThenBy(c => c.Id)
    .First();
\end{codeblock}

\unnumberedsubsection{D.3. Gộp nhóm sinh danh sách đi chợ}

{\raggedright Trích từ \codelong{src/CookingAdvisor/Services/ShoppingListService.cs}:\par}

\begin{codeblock}
var plan = await db.MenuPlans
    .Include(p => p.Items).ThenInclude(i => i.Recipe)
        .ThenInclude(r => r.RecipeIngredients)
    .FirstOrDefaultAsync(p =>
        p.Id == menuPlanId && p.UserId == userId);

if (plan is null)
    throw new InvalidOperationException(
        "Menu plan not found.");

var aggregated = plan.Items
    .SelectMany(i => i.Recipe.RecipeIngredients)
    .GroupBy(ri => (ri.IngredientId, ri.Unit))
    .Select(g => new ShoppingListItem
    {
        IngredientId = g.Key.IngredientId,
        Unit = g.Key.Unit,
        Quantity = g.Sum(ri => ri.Quantity)
    })
    .ToList();
\end{codeblock}

\unnumberedsection{Phụ lục E. Đặc tả chi tiết ba ca sử dụng lõi}

Ba ca sử dụng ứng với ba thuật toán lõi được đặc tả đầy đủ dưới đây; các ca còn
lại đã tóm tắt ở Bảng~\ref{tab:uc-tom-tat}.

\unnumberedsubsection{E.1. UC01. Gợi ý món theo nguyên liệu sẵn có}

\begin{longtable}{|P{3.4cm}|P{10.6cm}|}
\caption{Đặc tả ca sử dụng UC01. Gợi ý món theo nguyên liệu sẵn có}
\label{tab:uc01}\\
\hline
\bfseries Thành phần & \bfseries Nội dung \\ \hline
\endfirsthead
\hline
\bfseries Thành phần & \bfseries Nội dung \\ \hline
\endhead
\endfoot
Tác nhân & Khách, Người dùng, Quản trị viên \\ \hline
Mục tiêu & Tìm những món có thể nấu từ tập nguyên liệu đang có, xếp theo mức độ
  phù hợp \\ \hline
Tiền điều kiện & Kho dữ liệu đã có món ăn và nguyên liệu; không cần đăng nhập \\ \hline
Luồng chính & Người dùng mở mục Gợi ý nguyên liệu; hệ thống nạp danh sách nguyên
  liệu sắp theo tên; người dùng tick các nguyên liệu đang có rồi bấm Gợi ý món;
  hệ thống tính độ phủ cho từng món, xếp hạng theo bốn tiêu chí và hiển thị lưới
  kết quả kèm nhãn tình trạng nguyên liệu \\ \hline
Luồng thay thế & Nếu người dùng chưa chọn nguyên liệu nào, hệ thống giữ nguyên
  trạng thái ban đầu kèm hướng dẫn thay vì trả về kết quả rỗng. Nếu không có món
  nào chung nguyên liệu với tập đã chọn, hệ thống hiển thị trạng thái rỗng \\ \hline
Hậu điều kiện & Danh sách món được hiển thị theo đúng thứ tự ưu tiên đã thiết
  kế; hệ thống không ghi dữ liệu nào xuống cơ sở dữ liệu \\ \hline
\end{longtable}

\unnumberedsubsection{E.2. UC02. Sinh thực đơn tuần tự động}

\begin{longtable}{|P{3.4cm}|P{10.6cm}|}
\caption{Đặc tả ca sử dụng UC02. Sinh thực đơn tuần tự động}
\label{tab:uc02}\\
\hline
\bfseries Thành phần & \bfseries Nội dung \\ \hline
\endfirsthead
\hline
\bfseries Thành phần & \bfseries Nội dung \\ \hline
\endhead
\endfoot
Tác nhân & Người dùng \\ \hline
Mục tiêu & Sinh tự động một thực đơn đủ 21 suất ăn cho bảy ngày \\ \hline
Tiền điều kiện & Người dùng đã đăng nhập; kho dữ liệu có ít nhất một món thỏa bộ
  lọc vùng miền \\ \hline
Luồng chính & Người dùng mở mục Thực đơn tuần, nhập tên thực đơn, chọn tuần bắt
  đầu và vùng miền nếu muốn, rồi bấm Tạo thực đơn; hệ thống kiểm tra dữ liệu
  nhập, gọi thuật toán sinh 21 suất ăn, lưu thực đơn cùng 21 suất và chuyển sang
  trang chi tiết \\ \hline
Luồng thay thế & Nếu dữ liệu nhập không hợp lệ, hệ thống hiển thị lỗi ngay dưới
  ô nhập tương ứng và giữ nguyên giá trị đã nhập. Nếu không có món nào thỏa bộ
  lọc vùng miền, thuật toán ném ngoại lệ và hệ thống hiển thị thông báo không
  sinh được thực đơn \\ \hline
Hậu điều kiện & Một bản ghi \code{MenuPlan} và 21 bản ghi \code{MenuPlanItem}
  được lưu, gắn với đúng tài khoản đang đăng nhập \\ \hline
\end{longtable}

\unnumberedsubsection{E.3. UC04. Tạo danh sách đi chợ}

\begin{longtable}{|P{3.4cm}|P{10.6cm}|}
\caption{Đặc tả ca sử dụng UC04. Tạo danh sách đi chợ}
\label{tab:uc04}\\
\hline
\bfseries Thành phần & \bfseries Nội dung \\ \hline
\endfirsthead
\hline
\bfseries Thành phần & \bfseries Nội dung \\ \hline
\endhead
\endfoot
Tác nhân & Người dùng \\ \hline
Mục tiêu & Tổng hợp toàn bộ nguyên liệu của thực đơn thành danh sách cần mua \\ \hline
Tiền điều kiện & Thực đơn đã tồn tại và thuộc về người dùng đang đăng nhập \\ \hline
Luồng chính & Người dùng bấm Tạo danh sách ở trang chi tiết thực đơn; hệ thống
  nạp thực đơn kèm điều kiện lọc theo tài khoản, duyệt toàn bộ suất ăn, gộp
  nhóm theo cặp nguyên liệu và đơn vị, cộng dồn khối lượng rồi lưu danh sách và
  chuyển sang trang chi tiết danh sách \\ \hline
Luồng thay thế & Nếu thực đơn không tồn tại hoặc không thuộc về người dùng, hệ
  thống trả về lỗi không tìm thấy. Nếu thực đơn đã có danh sách đi chợ, hệ thống
  xóa toàn bộ dòng cũ rồi ghi lại thay vì tạo bản ghi mới \\ \hline
Hậu điều kiện & Một bản ghi \code{ShoppingList} và các bản ghi
  \code{ShoppingListItem} đã gộp được lưu \\ \hline
\end{longtable}

\unnumberedsection{Phụ lục F. Danh sách ca kiểm thử}

\begin{longtable}{|P{1.3cm}|P{5.4cm}|P{5.2cm}|P{1.2cm}|}
\caption{Bảng ca kiểm thử của bốn nhóm chức năng}
\label{tab:tc}\\
\hline
\bfseries Mã & \bfseries Tên ca kiểm thử & \bfseries Mục tiêu kiểm chứng &
\bfseries Kết quả \\ \hline
\endfirsthead
\hline
\bfseries Mã & \bfseries Tên ca kiểm thử & \bfseries Mục tiêu kiểm chứng &
\bfseries Kết quả \\ \hline
\endhead
\endfoot
\multicolumn{4}{|P{14.37cm}|}{\bfseries Nhóm 1. Thuật toán gợi ý theo nguyên liệu (SuggestionService)} \\ \hline
TC-01 & \codelong{SuggestAsync\_\allowbreak{}EmptyIngredientSet\_\allowbreak{}ReturnsEmpty} & Tập nguyên liệu
  rỗng trả về danh sách rỗng & Đạt \\ \hline
TC-02 & \codelong{SuggestAsync\_\allowbreak{}AllIngredientsOwned\_\allowbreak{}IsCookableWithNoMissing}
  & Có đủ nguyên liệu thì món được đánh dấu nấu được ngay, danh sách thiếu
  rỗng & Đạt \\ \hline
TC-03 & \codelong{SuggestAsync\_\allowbreak{}SomeIngredientsMissing\_\allowbreak{}ListsMissingNamesAndCoverage}
  & Liệt kê đúng tên nguyên liệu thiếu và tính đúng độ phủ & Đạt \\ \hline
TC-04 & \codelong{SuggestAsync\_\allowbreak{}NoMatchedIngredient\_\allowbreak{}RecipeIsExcluded} & Món
  không chung nguyên liệu nào bị loại khỏi kết quả & Đạt \\ \hline
TC-05 & \codelong{SuggestAsync\_\allowbreak{}RanksCookableFirstThenCoverageThenFewerMissing} &
  Thứ tự xếp hạng đúng theo bốn tiêu chí & Đạt \\ \hline
TC-06 &
  \codelong{GetIngredientOptionsAsync\_\allowbreak{}ReturnsAllIngredientsOrderedByName} & Danh
  sách nguyên liệu trả về đầy đủ và sắp theo tên & Đạt \\ \hline
\multicolumn{4}{|P{14.37cm}|}{\bfseries Nhóm 2. Thuật toán sinh thực đơn tuần (MenuPlannerService)} \\ \hline
TC-07 & \codelong{GenerateWeeklyPlanAsync\_\allowbreak{}FillsAll21Slots} & Sinh đủ 21 suất ăn &
  Đạt \\ \hline
TC-08 & \codelong{GenerateWeeklyPlanAsync\_\allowbreak{}EnoughRecipes\_\allowbreak{}NoRepeatsWithinWeek}
  & Thư viện đủ lớn thì không lặp món trong tuần & Đạt \\ \hline
TC-09 &
  \codelong{GenerateWeeklyPlanAsync\_\allowbreak{}FewRecipes\_\allowbreak{}AllowsRepeatsButStillFills21} & Thư
  viện nhỏ vẫn lấp đủ 21 suất bằng cách cho lặp & Đạt \\ \hline
TC-10 & \codelong{GenerateWeeklyPlanAsync\_\allowbreak{}FewRecipes\_\allowbreak{}SpreadsRepeatsEvenly} &
  Khi buộc lặp thì rải đều, không dồn vào một món & Đạt \\ \hline
TC-11 &
  \codelong{GenerateWeeklyPlanAsync\_\allowbreak{}ScarceBreakfastRecipes\_\allowbreak{}VariesBreakfastAcrossWeek}
  & Bữa sáng thiếu món vẫn đa dạng qua các ngày & Đạt \\ \hline
TC-12 & \codelong{GenerateWeeklyPlanAsync\_\allowbreak{}RespectsRegionFilter} & Ràng buộc
  vùng miền không bị nới & Đạt \\ \hline
TC-13 & \codelong{GenerateWeeklyPlanAsync\_\allowbreak{}PrefersFavoriteRecipe} & Món yêu thích
  được ưu tiên chọn & Đạt \\ \hline
TC-14 & \codelong{GenerateWeeklyPlanAsync\_\allowbreak{}NoRecipesAvailable\_\allowbreak{}Throws} & Không
  có món nào thì ném ngoại lệ & Đạt \\ \hline
TC-15 & \codelong{GenerateWeeklyPlanAsync\_\allowbreak{}PersistsMenuPlanAndItems} & Thực đơn và
  các suất được lưu xuống cơ sở dữ liệu & Đạt \\ \hline
\multicolumn{4}{|P{14.37cm}|}{\bfseries Nhóm 3. Thuật toán sinh danh sách đi chợ (ShoppingListService)} \\ \hline
TC-16 &
  \codelong{GenerateFromMenuPlanAsync\_\allowbreak{}AggregatesSameIngredientAndUnitAcrossRecipes}
  & Cùng nguyên liệu, cùng đơn vị ở nhiều món được cộng dồn & Đạt \\ \hline
TC-17 &
  \codelong{GenerateFromMenuPlanAsync\_\allowbreak{}KeepsSeparateRowsForDifferentUnits} & Cùng
  nguyên liệu khác đơn vị được tách thành hai dòng & Đạt \\ \hline
TC-18 &
  \codelong{GenerateFromMenuPlanAsync\_\allowbreak{}RepeatedRecipeOccurrences\_\allowbreak{}MultipliesQuantity}
  & Món lặp trong tuần thì khối lượng nhân lên tương ứng & Đạt \\ \hline
TC-19 & \codelong{GenerateFromMenuPlanAsync\_\allowbreak{}Regenerate\_\allowbreak{}ReplacesOldItems} &
  Tạo lại thì ghi đè, không nhân đôi dữ liệu & Đạt \\ \hline
TC-20 &
  \codelong{GenerateFromMenuPlanAsync\_\allowbreak{}EmptyPlan\_\allowbreak{}CreatesShoppingListWithNoItems} &
  Thực đơn rỗng vẫn tạo được danh sách rỗng hợp lệ & Đạt \\ \hline
TC-21 & \codelong{GenerateFromMenuPlanAsync\_\allowbreak{}NonexistentPlan\_\allowbreak{}Throws} & Thực
  đơn không tồn tại thì ném ngoại lệ & Đạt \\ \hline
TC-22 & \codelong{GenerateFromMenuPlanAsync\_\allowbreak{}PlanOwnedByDifferentUser\_\allowbreak{}Throws} &
  Không truy cập được thực đơn của người dùng khác & Đạt \\ \hline
TC-23 & \codelong{GenerateFromMenuPlanAsync\_\allowbreak{}PersistsMenuPlanIdAndUserId} &
  Lưu đúng liên kết tới thực đơn và người dùng & Đạt \\ \hline
\multicolumn{4}{|P{14.37cm}|}{\bfseries Nhóm 4. Chức năng tìm kiếm và sắp xếp (RecipeService)} \\ \hline
TC-24 & \codelong{Defaults\_\allowbreak{}to\_\allowbreak{}alphabetical\_\allowbreak{}order} & Mặc định sắp theo tên món &
  Đạt \\ \hline
TC-25 & \codelong{Sorts\_\allowbreak{}by\_\allowbreak{}total\_\allowbreak{}time\_\allowbreak{}using\_\allowbreak{}prep\_\allowbreak{}plus\_\allowbreak{}cook} & Sắp theo
  tổng thời gian chuẩn bị cộng nấu & Đạt \\ \hline
TC-26 & \codelong{Sorts\_\allowbreak{}by\_\allowbreak{}calories\_\allowbreak{}in\_\allowbreak{}both\_\allowbreak{}directions} (calo tăng dần) & Sắp
  đúng chiều tăng & Đạt \\ \hline
TC-27 & \codelong{Sorts\_\allowbreak{}by\_\allowbreak{}calories\_\allowbreak{}in\_\allowbreak{}both\_\allowbreak{}directions} (calo giảm dần) &
  Sắp đúng chiều giảm & Đạt \\ \hline
TC-28 & \codelong{Breaks\_\allowbreak{}ties\_\allowbreak{}by\_\allowbreak{}name\_\allowbreak{}so\_\allowbreak{}paging\_\allowbreak{}cannot\_\allowbreak{}repeat\_\allowbreak{}a\_\allowbreak{}dish} & Có
  khóa phá vỡ thế cân bằng, phân trang không lặp món & Đạt \\ \hline
TC-29 & \codelong{Sort\_\allowbreak{}survives\_\allowbreak{}alongside\_\allowbreak{}a\_\allowbreak{}filter\_\allowbreak{}and\_\allowbreak{}is\_\allowbreak{}echoed\_\allowbreak{}back}
  & Sắp xếp hoạt động cùng bộ lọc và được phản hồi lại giao diện & Đạt \\ \hline
\end{longtable}

\unnumberedsection{Phụ lục G. Cấu trúc chi tiết các bảng cơ sở dữ liệu}

Kiểu dữ liệu ghi trong bảng là kiểu trên SQL Server sau khi EF Core ánh xạ từ
lớp thực thể; các cột khóa chính tự tăng (\code{Id}, kiểu int IDENTITY) không
lặp lại trong phần mô tả. Vai trò và khóa của từng bảng đã tóm tắt ở
Bảng~\ref{tab:tom-tat-bang}.

\begin{longtable}{|P{4.2cm}|P{3.0cm}|P{6.2cm}|}
\caption{Cấu trúc chi tiết các bảng của cơ sở dữ liệu}
\label{tab:cau-truc-bang}\\
\hline
\bfseries Trường & \bfseries Kiểu & \bfseries Mô tả \\ \hline
\endfirsthead
\hline
\bfseries Trường & \bfseries Kiểu & \bfseries Mô tả \\ \hline
\endhead
\endfoot
\multicolumn{3}{|P{14.24cm}|}{\bfseries Bảng Categories: danh mục món ăn, làm bộ lọc
  trên trang danh sách} \\ \hline
\code{Name} & nvarchar, NOT NULL & Tên danh mục, ví dụ Món canh, Món xào \\ \hline
\multicolumn{3}{|P{14.24cm}|}{\bfseries Bảng Ingredients: kho nguyên liệu dùng chung} \\ \hline
\code{Name} & nvarchar, NOT NULL & Tên nguyên liệu \\ \hline
\code{Unit} & nvarchar, NOT NULL & Đơn vị đo mặc định khi hiển thị \\ \hline
\code{Group} & nvarchar, NOT NULL & Nhóm nguyên liệu (Rau củ, Thịt, Gia vị...),
  căn cứ gom nhóm của danh sách đi chợ \\ \hline
\multicolumn{3}{|P{14.24cm}|}{\bfseries Bảng Recipes: bảng trung tâm, lưu thông tin món
  ăn} \\ \hline
\code{Name} & nvarchar, NOT NULL & Tên món, bắt buộc \\ \hline
\code{Description} & nvarchar, NULL & Mô tả ngắn \\ \hline
\code{Instructions} & nvarchar, NULL & Các bước nấu, lưu dạng văn xuôi \\ \hline
\code{Servings} & int & Số khẩu phần \\ \hline
\code{PrepMinutes} & int & Thời gian chuẩn bị, tính bằng phút \\ \hline
\code{CookMinutes} & int & Thời gian nấu, tính bằng phút \\ \hline
\code{Difficulty} & int & Độ khó: 0 Dễ, 1 Trung bình, 2 Khó \\ \hline
\code{Region} & int & Vùng miền: 0 Bắc, 1 Trung, 2 Nam \\ \hline
\codelong{CaloriesPerServing} & int & Năng lượng mỗi khẩu phần, kcal \\ \hline
\code{ImageUrl} & nvarchar, NULL & Đường dẫn ảnh minh họa \\ \hline
\codelong{SuitableMealTypes} & int & Cờ nhị phân các bữa phù hợp \\ \hline
\code{CategoryId} & int, FK & Khóa ngoại tới \code{Categories} \\ \hline
\multicolumn{3}{|P{14.24cm}|}{\bfseries Bảng RecipeIngredients: bảng trung gian món ăn -
  nguyên liệu, mấu chốt của cả ba thuật toán} \\ \hline
\code{RecipeId} & int, PK, FK & Khóa chính thành phần, trỏ tới \code{Recipes} \\ \hline
\code{IngredientId} & int, PK, FK & Khóa chính thành phần, trỏ tới
  \code{Ingredients} \\ \hline
\code{Quantity} & decimal(10,2) & Khối lượng cần dùng \\ \hline
\code{Unit} & nvarchar, NOT NULL & Đơn vị đo dùng trong công thức này \\ \hline
\multicolumn{3}{|P{14.24cm}|}{\bfseries Bảng MenuPlans: thông tin chung của một thực đơn
  tuần} \\ \hline
\code{UserId} & nvarchar, FK & Khóa ngoại tới bảng người dùng của Identity \\ \hline
\code{Name} & nvarchar, NOT NULL & Tên thực đơn do người dùng đặt \\ \hline
\code{WeekStartDate} & date & Ngày bắt đầu của tuần được lập \\ \hline
\code{CreatedAt} & datetime2 & Thời điểm tạo thực đơn \\ \hline
\multicolumn{3}{|P{14.24cm}|}{\bfseries Bảng MenuPlanItems: từng ô của lưới thực đơn, một
  thực đơn đầy đủ gồm 21 bản ghi} \\ \hline
\code{MenuPlanId} & int, FK & Khóa ngoại tới \code{MenuPlans} \\ \hline
\code{DayOfWeek} & int & Thứ trong tuần, nhận giá trị từ 0 tới 6 \\ \hline
\code{MealType} & int & Bữa ăn: 0 Sáng, 1 Trưa, 2 Tối \\ \hline
\code{RecipeId} & int, FK & Khóa ngoại tới \code{Recipes} \\ \hline
\multicolumn{3}{|P{14.24cm}|}{\bfseries Bảng Favorites: dấu yêu thích của người dùng đối
  với một món} \\ \hline
\code{UserId} & nvarchar, PK, FK & Khóa chính thành phần, trỏ tới bảng người
  dùng \\ \hline
\code{RecipeId} & int, PK, FK & Khóa chính thành phần, trỏ tới \code{Recipes} \\ \hline
\multicolumn{3}{|P{14.24cm}|}{\bfseries Bảng ShoppingLists: thông tin chung của một danh
  sách đi chợ} \\ \hline
\code{MenuPlanId} & int, FK, UNIQUE & Khóa ngoại tới \code{MenuPlans}, có ràng
  buộc duy nhất \\ \hline
\code{UserId} & nvarchar, FK & Khóa ngoại tới bảng người dùng \\ \hline
\code{CreatedAt} & datetime2 & Thời điểm sinh danh sách \\ \hline
\multicolumn{3}{|P{14.24cm}|}{\bfseries Bảng ShoppingListItems: từng dòng nguyên liệu đã
  gộp} \\ \hline
\code{ShoppingListId} & int, FK & Khóa ngoại tới \code{ShoppingLists} \\ \hline
\code{IngredientId} & int, FK & Khóa ngoại tới \code{Ingredients} \\ \hline
\code{Quantity} & decimal(10,2) & Tổng khối lượng cần mua sau khi cộng dồn \\ \hline
\code{Unit} & nvarchar, NOT NULL & Đơn vị đo, là một nửa của khóa gộp nhóm \\ \hline
\code{IsPurchased} & bit & Đã mua hay chưa, phục vụ thanh tiến độ \\ \hline
\end{longtable}

\unnumberedsection{Phụ lục H. Mã giả ba thuật toán lõi}

Các quyết định cài đặt của ba thuật toán được phân tích ở
mục~\ref{sec:cai-dat-thuat-toan}, tham chiếu theo số dòng của mã giả dưới đây.

\unnumberedsubsection{H.1. Thuật toán gợi ý theo nguyên liệu}

\begin{codeblock}
THUẬT TOÁN GoiYTheoNguyenLieu
Đầu vào : A - tập mã nguyên liệu người dùng đang có
Đầu ra  : danh sách món đã xếp hạng kèm nguyên liệu thiếu

 1  nếu A rỗng thì
 2      trả về danh sách rỗng
 3  hết nếu
 4
 5  A <- chuyển A thành bảng băm      // tra cứu O(1)
 6  R <- truy vấn món có ít nhất 1 nguyên liệu thuộc A
 7  KetQua <- danh sách rỗng
 8
 9  với mỗi món r thuộc R lặp
10      I_r      <- tập nguyên liệu của r
11      missing  <- { i thuộc I_r | i không thuộc A }
12      missing  <- sắp xếp missing theo bảng chữ cái
13      matched  <- |I_r| - |missing|
14      coverage <- matched / |I_r|
15      canCook  <- (missing rỗng)
16      thêm (r, missing, coverage, canCook) vào KetQua
17  hết lặp
18
19  sắp xếp KetQua theo thứ tự từ điển:
20      khóa 1: canCook    giảm dần  // nấu được ngay trước
21      khóa 2: coverage   giảm dần
22      khóa 3: |missing|  tăng dần
23      khóa 4: tên món    tăng dần  // giữ thứ tự tất định
24
25  trả về KetQua
\end{codeblock}

\unnumberedsubsection{H.2. Thuật toán sinh thực đơn tuần}

\begin{codeblock}
THUẬT TOÁN SinhThucDonTuan
Đầu vào : userId, tenThucDon, tuanBatDau, vungMien (tuỳ chọn)
Đầu ra  : đối tượng MenuPlan chứa 21 MenuPlanItem
Hằng số : MUC_TIEU_CALO_NGAY = 2000
          BUA[] = [ (Sáng, cờ Breakfast), (Trưa, cờ Lunch),
                    (Tối, cờ Dinner) ]

 1  candidates <- các món thoả vùng miền   // ràng buộc CỨNG
 2  nếu candidates rỗng thì ném ngoại lệ
 3
 4  favoriteIds <- tập mã món yêu thích của userId
 5  used      <- tập rỗng      // món đã dùng trong tuần
 6  useCounts <- từ điển rỗng  // số lần mỗi món đã dùng
 7  items     <- danh sách rỗng
 8
 9  với d <- 0 đến 6 lặp                   // 7 ngày
10      dayCalories <- 0
11      với m <- 0 đến 2 lặp               // 3 bữa
12          (bua, co) <- BUA[m]
13          mealTarget <- MUC_TIEU_CALO_NGAY x (m+1) / 3
14
15          suitable <- { c thuộc candidates | c hợp cờ co }
16          pool     <- suitable \ used
17          repeating <- (pool rỗng)
18
19          nếu repeating: pool <- suitable   // nới KHÔNG LẶP
20          nếu pool rỗng: pool <- candidates // nới HỢP BỮA
21
22          nếu repeating thì
23              ordered <- sắp pool theo: useCounts[c] tăng,
24                            rồi (c thuộc favoriteIds) giảm
25          ngược lại
26              ordered <- sắp pool theo: c thuộc favoriteIds
27          hết nếu
28
29          chosen <- phần tử đầu của ordered sau khi sắp:
30                  |dayCalories + c.Calo - mealTarget| tăng,
31                    rồi c.Id tăng dần
32
33          dayCalories <- dayCalories + chosen.Calo
34          thêm chosen vào used
35          useCounts[chosen] <- useCounts[chosen] + 1
36          thêm MenuPlanItem(d, bua, chosen) vào items
37      hết lặp
38  hết lặp
39
40  lưu MenuPlan(userId, tenThucDon, tuanBatDau, items)
41  trả về MenuPlan
\end{codeblock}

\unnumberedsubsection{H.3. Thuật toán sinh danh sách đi chợ}

\begin{codeblock}
THUẬT TOÁN SinhDanhSachDiCho
Đầu vào : userId, menuPlanId
Đầu ra  : đối tượng ShoppingList

 1  plan <- nạp MenuPlan kèm Items, Recipe, RecipeIngredients
 2          với điều kiện Id = menuPlanId VÀ UserId = userId
 3  nếu plan là null thì ném ngoại lệ  // chặn truy cập chéo
 4
 5  tatCa <- danh sách rỗng
 6  với mỗi suất s thuộc plan.Items lặp  // duyệt theo SUẤT
 7      thêm toàn bộ s.Recipe.RecipeIngredients vào tatCa
 8  hết lặp
 9
10  nhom <- gộp nhóm tatCa theo khóa (IngredientId, Unit)
11  kq <- danh sách rỗng
12  với mỗi nhóm g thuộc nhom lặp
13      thêm ShoppingListItem(g.Key, tổng g.Quantity) vào kq
14  hết lặp
15
16  list <- tìm ShoppingList có MenuPlanId = menuPlanId
17  nếu list tồn tại thì
18      xóa toàn bộ dòng cũ của list      // ghi đè bản cũ
19  ngược lại
20      list <- tạo ShoppingList mới
21  hết nếu
22
23  gán kq vào list.Items và lưu
24  trả về list
\end{codeblock}
